% Chapter 2: Background

\chapter{Background}

This chapter reviews the background material required to understand the PVE-ZKP
system. The focus is on three topics: zero-knowledge proof systems and the
Groth16 construction, tooling for building and verifying circuits in practice,
and conventional web input validation patterns and their limitations.

\section{Zero-Knowledge Proofs and Groth16}

Zero-knowledge proofs allow one party (the prover) to convince another party
(the verifier) that a statement is true without revealing the underlying
witness. Modern zk-SNARK constructions, such as Groth16, provide succinct,
non-interactive proofs and fast verification, which makes them attractive for
use in deployed systems.

The Groth16 protocol represents a computation as a rank-1 constraint system
(R1CS) and, after a trusted setup phase, allows the prover to generate a proof
that they know a witness satisfying all constraints. Verification uses a small
set of group operations and pairing checks, with complexity independent of the
witness size. In PVE-ZKP, we rely on existing implementations of Groth16 in the
snarkjs library rather than re-deriving the protocol.

\section{Circom, snarkjs, and circomlib}

Circom is a domain-specific language for writing arithmetic circuits that can
be compiled to R1CS and used with proving systems such as Groth16. Circom
circuits define signals and constraints in a modular way, making it easier to
encode typical validation patterns such as range checks, set membership, and
string-format constraints.

The snarkjs toolchain provides command-line and JavaScript APIs for compiling
Circom circuits, running trusted setup ceremonies, generating proofs, and
verifying proofs. The circomlib and circomlibjs libraries supply reusable
components such as hash functions (e.g., Poseidon) and field-friendly
primitives. PVE-ZKP uses Circom 2.1 and snarkjs in its build scripts and
runtime services.

\section{Conventional Web Input Validation}

In conventional web architectures, input validation is performed in two places:
client-side checks in the browser (for responsiveness and user experience) and
server-side checks in the receiving application backend. Client-side checks are
easily bypassed by an adversary who can craft arbitrary HTTP requests, so
security ultimately relies on the server-side logic.

Third-party validation services have emerged to provide centralized policy
engines for common checks such as email verification, address normalization, and
fraud detection. These services typically require the application backend to
send raw user data to the validator and then accept or reject submissions based
on the returned decision. While convenient, this approach increases the number
of parties that see sensitive data and creates additional trust assumptions.

\section{Related Work on ZK Circuit Analysis and Secure Aggregation}

The PVE-ZKP system is complementary to two lines of work. First, security
analysis tools for ZK circuits, such as those that search for underconstrained
signals or unsafe compositions, help developers detect mistakes that could
invalidate soundness guarantees. These tools operate at the circuit-design
level and can be integrated into the PVE-ZKP development workflow but do not
address web input validation architecture directly.

Second, schemes for secure aggregation and privacy-preserving federated
learning use ZKPs or related primitives to ensure that individual client
contributions satisfy policy constraints without revealing the contributions
themselves. PVE-ZKP instead targets single-form submissions in classical web
applications, but shares the goal of separating validation from disclosure.

This background motivates the design of PVE-ZKP as an application of existing
ZK tools and ideas to a different part of the web security landscape: input
validation for browser-based forms.
